\documentclass[10pt,a4paper]{amsart}
\usepackage[margin=19mm]{geometry}
\usepackage{amsmath,amssymb,mathtools}
\usepackage[scr=rsfs,cal=euler]{mathalfa}
\usepackage{xfrac}
\usepackage[unicode,hidelinks,bookmarks=false]{hyperref}
\newcommand{\ie}{i.e.\ }
\newcommand{\cf}{cf.\ }
\newcommand{\resp}{respectively\ }
\newcommand{\Subsection}{\subsection}
\providecommand{\nobreakdash}{\nobreak\hbox{-}}
\setlength{\emergencystretch}{2em}
\multlinegap0pt
\usepackage{amssymb}
\usepackage{enumerate}
\usepackage{dsfont}
\newtheorem{definition}{Definition}
\newtheorem{theorem}{Theorem}
\newtheorem{lemma}{Lemma}
\newtheorem{proposition}{Proposition}
\def \R {\mathbb{R}}
\newcommand{\da}{\rangle}
\newcommand{\la}{\lambda}
\newcommand{\p}{\partial}
\newcommand{\sa}{\langle}
\newcommand{\tr}{\operatorname{tr}}
\def \vf{\varphi}
\title{Strichartz estimates for a class of Schr\"odinger equations with a drift}
\author{Federico Buseghin, Nicola Garofalo}
\date{Source: arXiv:2503.12606v4 (CC BY 4.0)}
\begin{document}
\maketitle

\noindent\emph{Source and licence.} Strichartz estimates for a class of Schr\"odinger equations with a drift. Version arXiv:2503.12606v4 (CC BY 4.0), distributed by arXiv under the Creative Commons Attribution 4.0 International licence, http://creativecommons.org/licenses/by/4.0/, as recorded in the arXiv licence metadata for this exact version and shown on the abstract page https://arxiv.org/abs/2503.12606v4. Full licence notice: \texttt{licenses/arxiv-2503.12606-CC-BY-4.0.txt}.\par\smallskip

\noindent\textit{Editorial scope.} Counted unit: the article's theorem T:trB (source0.tex lines 592--616). The article's complete original proof of it is delivered with it (source0.tex lines 1704--1753, one \texttt{proof} environment of 50 lines, which the article places away from the statement). No mathematics is written, repaired or re-derived: the statement and the proof are the article's own lines, in the article's own notation, and no retained line is substituted or re-flowed.  Retained uncounted as the source-local context the proof needs: lines 351--353; lines 442--448; lines 522--524; lines 1562--1567; lines 1583--1588; lines 2469--2471; lines 2473--2475; lines 2502--2504; lines 2514--2516; lines 2520--2527.  Nothing was omitted from any retained span, so the package carries no omission record.  No retained line contains a citation, so the wrapper carries no bibliography and no bibliography entry.  No retained line contains a figure environment, an image-inclusion command or drawing code, so no image is delivered and the package declares no asset dependency.  Editorial formatting only, in addition: an AMS \texttt{amsart} preamble that restates the article's own declarations and macros used by the retained lines, namely the environments \texttt{definition}, \texttt{theorem}, \texttt{lemma}, \texttt{proposition} on separate counters, together with the standard packages those lines need.  The retained environments print in this document as Definition 1, Theorem 1, Lemma 1, Proposition 1 in order of appearance, whereas the article prints the same items with the numbering of its own full document.  The complete native source of the article as distributed in the arXiv e-print for this exact version is bundled unchanged as \texttt{sources/174866/source0.tex}.
\begin{equation}\label{A}
\p_t u - i \tr(Q\nabla^2 u) -  \sa Bx,\nabla u\da   = F(x,t),\ \ \ \ \ \ \ u(x,0) = \vf(x).
\end{equation}

\begin{definition}\label{D:hd}
We call the number  
\begin{equation}\label{hd}
D = p_0+3p_1+\ldots+(2r+1)p_r
\end{equation}
the \emph{local homogeneous dimension} of the Schr\"odinger operator \eqref{A}. 
\end{definition}

\begin{equation}\label{H0}
V(t) \ge \gamma\ t^{D} e^{t \tr B}.
\end{equation}

\begin{lemma}\label{L:ber}
The function $t\to \frac{V(t)}{e^{2 t \tr B}}$ is non-decreasing on $(0,\infty)$ and for $t\ge 1$ we have  
\begin{equation}\label{below}
\frac{V(t)}{e^{2 t \tr B}} \ge \tilde V(1)>0.
\end{equation}
\end{lemma}

\begin{proposition}\label{P:ii}
Suppose that $\operatorname{Rank}(Q) = n$, and that $\sigma(B)\subset i \mathbb{R}$. There exist $D_{\infty}\ge n$ and $\gamma>0$, such that for $t\ge 1$ one has
\begin{align*}
V(t)\ge \gamma t^{D_{\infty}}.
\end{align*}
\end{proposition}

\begin{equation}\label{quasi}
\max\{\Re(\lambda)\mid \lambda\in \sigma(B)\}< 0.
\end{equation}

\begin{equation}\label{quasinfty}
V(t) = \det Q(t) \ \longrightarrow\ V_\infty = \det Q_\infty > 0.
\end{equation}

\begin{equation}\label{voldil}
\overline V(t) = \gamma_0\ t^{D},\ \ \ \ \ t>0,
\end{equation}

\begin{equation}\label{vicinoa0}
V(t) \ge \gamma t^D,\ \ \ \ \ \ 0<t\le 1.
\end{equation}

\begin{proposition}\label{P:boom0}
The following is true:
\begin{itemize}
\item[(i)] If $\tr B \ge 0$, then there exists a constant $c_1>0$ such that $V(t) \geq c_1 t^2$ for all $t\geq 1$.
\item[(ii)] If $\tr B \ge 0$ and furthermore $\max\{\Re(\lambda)\mid \lambda\in \sigma(B)\}=L_0>0$, then there exists a constant $c_0$ such that $V(t)\geq c_0 e^{2 L_0 t}$ for all $t\geq 1$.
\item[(iii)] Finally, if $\max\{\Re(\lambda)\mid \lambda\in \sigma(B)\}\geq 0$, then there exists a constant $c_0>0$ such that $V(t)\geq c_0 t$, for every $t\ge 1$.
\end{itemize}
\end{proposition}

\begin{theorem}\label{T:trB}
Suppose that $Q$ and $B$ satisfy \emph{one} of the following hypothesis \emph{(i)-(iii)}:
\begin{itemize}
\item[(i)] $\sigma(B)\not\subset i\R$;
\item[(ii)] $\sigma\{B\}\subset i \R$ and $\operatorname{Rank}(Q) = n$; 
\item[(iii)] $B$ is nilpotent and of the form 
\begin{equation}\label{barB}
\overline B = \begin{pmatrix} 0 & 0 & \cdot & \cdot & \cdot & 0 & 0
\\
B_1 & 0 & \cdot & \cdot & \cdot & 0 & 0
\\
0 & B_2 & \cdot & \cdot & \cdot &  0 & 0
\\
\cdot & \cdot & \cdot & \cdot & \cdot & \cdot & \cdot
\\
\cdot & \cdot & \cdot & \cdot & \cdot & \cdot & \cdot
\\
\cdot & \cdot & \cdot & \cdot & \cdot & \cdot & \cdot
\\
0 & 0 & \cdot & \cdot & \cdot &  B_r & 0
\end{pmatrix}.
\end{equation}
\end{itemize}
Then \eqref{H0} in \emph{Hypothesis (A)} holds.
\end{theorem}

\begin{proof}[Proof of Theorem \ref{T:trB}]

We begin by noting that, based on \eqref{vicinoa0}, the estimate \eqref{H0} holds for $0<t\le 1$, regardless of the value of $\tr B$. We are thus left with proving that it is valid also for $t\ge 1$.
Assume now the hypothesis (i) in the statement of the theorem. There are three possible cases.

\medskip

\noindent \underline{Case (1)}: $\tr B<0$. If \eqref{quasi} holds, then for $t\ge 1$ the desired conclusion \eqref{H0} immediately follows  from \eqref{quasinfty} and the fact that $\tr B<0$. If instead there exists at least one $\la_0\in \sigma(B)$ such that $\Re \la_0\ge 0$, then by (iii) in Proposition \ref{P:boom0} we know that for some $c_0>0$ we have for $t\ge 1$
\[
V(t) \ge c_0 t.
\]
From this estimate we infer that the inequality \eqref{H0} holds for $t\ge 1$, provided that for some $C>0$
\[
e^{-t \tr B} \ge C t^{D-1}.
\]
Since $\tr B<0$, this is of course true.

\medskip

\noindent \underline{Case (2)}: $\tr B>0$. By appealing to \eqref{below} in Lemma \ref{L:ber}, we obtain for $t\ge 1$ 
\[
V(t) \ge \gamma_0\ e^{2 t \tr B},
\]
for some constant $\gamma_0>0$. Since we are assuming $\tr B>0$, it is clear from this estimate that, for some $\gamma>0$, the inequality \eqref{H0} holds true for every $t\ge 1$.

\medskip

\noindent \underline{Case (3)}: $\tr B = 0$. To complete the proof, we thus need to show that there exists $\gamma >0$ such that
\begin{equation}\label{op}
V(t) \ge \gamma\ t^D,\ \ \ \ \ \ \ t\ge 1.
\end{equation}
By the hypothesis $\sigma(B)\not\subset i \R$, there exists at least one $\la _0\in \sigma(B)$, such that $\Re \la_0 \not=0$. If $\Re \la_0>0$, then from (ii) in Proposition \ref{P:boom0} we have $V(t) \ge c_0 e^{2 \Re \la_0 t}$ for $t\ge 1$, and thus \eqref{op} follows. If instead $\Re \la_0<0$, then we claim that there must be at least one other $\la_1\in \sigma(B)$ such that $\Re \la_1>0$. Otherwise, if all other eigenvalues of $B$ had real part $\le 0$, we would have $\tr B < 0$, a contradiction. Again from (ii) in Proposition \ref{P:boom0} we have $V(t) \ge c_1 e^{2 \Re \la_1 t}$ for $t\ge 1$, and \eqref{op} follows. We have thus proved the theorem under the hypothesis (i).

Next, assume that hypothesis (ii) holds. Since $Q$ is invertible, by (b) following Definition \ref{D:hd}, we have $D = n$, and therefore \eqref{vicinoa0} implies
\[
V(t)\ge \gamma\ t^{n},\ \ \ \ \ 0\le t\le 1.
\]
Now, applying Proposition \ref{P:ii} we conclude that, for some $D_\infty \ge n$, we have
\begin{align*}
V(t)\ge \gamma\ t^{D_{\infty}},\ \ \ \ \ t\ge 1.
\end{align*} 
Combining the latter two inequalities, we reach the conclusion that \eqref{H0} holds for every $t\ge 0$.

Finally, assume that the hypothesis (iii) holds. In this case, the identity \eqref{voldil} gives
\[
V(t) = \gamma\ t^D,\ \ \ \ \ \ t>0,
\]
for some $\gamma >0$. Since, as we have noted, we must necessarily have $\tr B = 0$, we infer that \eqref{H0} holds in this case as well.

\end{proof}

\end{document}
